\chapter{Stochastic Bühlmann-Straub credibility}

Bühlmann and Straub \cite{BuhlmannStraub1970}, Sections 3, 5, and 6, model a loss ratio $X_l$
using a latent risk parameter $\Theta$. This chapter fixes the contract
conditions. Conditional on $\Theta$, the loss
ratios have common mean $M=\mu(\Theta)$ and process variance
$\sigma^2(\Theta)/P_l$, where $P_l$ is exposure. The structural process
parameter is $v=E[\sigma^2(\Theta)]$. Their Section 6 minimizes expected squared prediction
error. This chapter proves the one-risk moment calculation and obtains the
credibility factor in display (4), page 123. The scanned primary source is
available from E-Periodica under DOI \texttt{10.5169/seals-967024}.

\begin{definition}[Conditional moment model]\label{def:bsStochastic}\lean{VerifiedReserving.buhlmannStraubLatentSigma, VerifiedReserving.buhlmannStraubRiskMean, VerifiedReserving.buhlmannStraubRiskVariance, VerifiedReserving.buhlmannStraubResidual, VerifiedReserving.BuhlmannStraubMomentModel}\leanok
Let $\mathcal G=\sigma(\Theta)$, $S=\sigma^2(\Theta)$, and $e_l=X_l-M$.
For the finite exposure
sample, the model records
\[
  E[X_l\mid\mathcal G]=M,\qquad
  E[e_l^2\mid\mathcal G]=\frac S{P_l},\qquad
  E[e_le_r\mid\mathcal G]=0\quad(l\ne r),
\]
together with $E[S]=v$, $E[M]=m$, $E[(M-m)^2]=w$, positive exposures,
measurability, and the needed integrability. The source assumes conditional independence; the
Lean proof records only the conditional cross moments it consumes.
\end{definition}

\begin{definition}[Predictor, expected loss, and coefficient loss]\label{def:bsRisk}\lean{VerifiedReserving.buhlmannStraubWeightedMeanRv, VerifiedReserving.buhlmannStraubWeightedResidualTotal, VerifiedReserving.buhlmannStraubPredictor, VerifiedReserving.buhlmannStraubExpectedSqLoss, VerifiedReserving.buhlmannStraubRiskQuadratic}\leanok
Write $P=\sum_lP_l$, $\bar X=\sum_lP_lX_l/P$, and
\[
  \widehat M_z=m+z(\bar X-m),\qquad
  R(z)=E[(M-\widehat M_z)^2].
\]
The exposure-scaled coefficient loss is
\[
  Q(z)=Pw(1-z)^2+vz^2.
\]
This stochastic objective is separate from the deterministic penalized sample
loss in Definition~\ref{def:bs}.
\end{definition}

\begin{theorem}[Conditional residual moments]\label{thm:bsResidualMoments}\lean{VerifiedReserving.buhlmannStraubWeightedMeanRv_sub_riskMean, VerifiedReserving.condExp_buhlmannStraubResidual_eq_zero, VerifiedReserving.condExp_buhlmannStraubWeightedResidualTotal_eq_zero, VerifiedReserving.integral_buhlmannStraubWeightedResidualTotal_eq_zero, VerifiedReserving.integral_buhlmannStraubPredictor_eq_collectiveMean, VerifiedReserving.condExp_sq_buhlmannStraubWeightedResidualTotal, VerifiedReserving.integral_sq_buhlmannStraubWeightedMeanRv_sub_riskMean, VerifiedReserving.integral_buhlmannStraubRiskMean_mul_weightedResidual_eq_zero}\leanok
\uses{def:bsStochastic, def:bsRisk}
The weighted residual total has conditional mean zero and conditional second
moment $SP$. When the total exposure $P$ is nonzero, it follows that
\[
 E[(\bar X-M)^2]=\frac vP,
 \qquad E[(M-m)(\bar X-M)]=0.
\]
Every predictor $m+z(\bar X-m)$ in the scalar family has expectation $m$,
so it satisfies the source's expectation-unbiasedness constraint.
\end{theorem}
\begin{proof}\leanok Scale each residual by its exposure. Conditional cross
moments remove the off-diagonal terms, and
$P_l^2(S/P_l)=SP_l$. Pulling the centered latent mean through conditional
expectation removes the remaining cross term.\end{proof}

\begin{theorem}[Expected prediction loss]\label{thm:bsExpectedLoss}\lean{VerifiedReserving.buhlmannStraubExpectedSqLoss_eq, VerifiedReserving.buhlmannStraubExpectedSqLoss_scaled_eq_riskQuadratic}\leanok
\uses{thm:bsResidualMoments}
For every scalar coefficient $z$ and nonzero total exposure,
\[
  R(z)=(1-z)^2w+z^2\frac vP,
  \qquad PR(z)=Q(z).
\]
Thus the deterministic coefficient quadratic is derived from the expected
squared prediction loss, rather than postulated.
\end{theorem}
\begin{proof}\leanok Expand
$M-\widehat M_z=(1-z)(M-m)-z(\bar X-M)$ and apply
Theorem~\ref{thm:bsResidualMoments}.\end{proof}

\begin{theorem}[Credibility is the risk minimizer]\label{thm:bsStochasticMinimizer}\lean{VerifiedReserving.buhlmannStraubRiskQuadratic_eq_at_credibility_add, VerifiedReserving.buhlmannStraubCredibility_minimizes_riskQuadratic, VerifiedReserving.buhlmannStraubCredibility_minimizes_expectedSqLoss, VerifiedReserving.buhlmannStraubPredictor_credibility_eq_estimate}\leanok
\uses{thm:bsExpectedLoss, thm:bsCredibility}
If $v+Pw>0$, then
\[
 Q(z)=Q(Z)+(v+Pw)(z-Z)^2,
 \qquad Z=\frac{Pw}{v+Pw}.
\]
For $P>0$, $Z$ minimizes the actual expected squared prediction loss. The
predictor $m+Z(\bar X-m)$ is pointwise equal to
\texttt{buhlmannStraubEstimate}. Supplying the source's empirical portfolio
mean as $m$ gives its display (4) estimator.
\end{theorem}
\begin{proof}\leanok Complete the square in $Q$, use positivity to compare
scaled expected losses, and apply the deterministic credibility-form identity.\end{proof}

\begin{theorem}[Finite process-shock witness]\label{thm:bsWitness}\lean{VerifiedReserving.BuhlmannStraubWitness.observation_nondegenerate, VerifiedReserving.BuhlmannStraubWitness.momentModel, VerifiedReserving.BuhlmannStraubWitness.exists_process_nondegenerate_buhlmannStraubModel}\leanok
\uses{def:bsStochastic}
There is a two-outcome probability model with one unit exposure, constant
latent mean $1$, and observed loss ratio $0$ or $2$ with equal probability.
It satisfies the moment model with process variance $v=1>0$ and has a
nonconstant observation. Its between-risk variance is zero; the general
theorems do not require that specialization.
\end{theorem}
\begin{proof}\leanok Conditional expectation given the trivial
$\sigma$-algebra is the two-point average.\end{proof}
